% Generated draft: requires mathematical and source review.
\documentclass[11pt]{article}
\usepackage[margin=25mm]{geometry}
\usepackage{fontspec}
\setmainfont{Latin Modern Roman}
\usepackage{amsmath,amssymb,mathtools,unicode-math}
\setmathfont{Latin Modern Math}
\usepackage{enumitem,needspace,xurl,hyperref}
\setlength{\parindent}{0pt}
\setlength{\parskip}{5pt}
\setlength{\emergencystretch}{3em}
\newcommand{\N}{\mathbb{N}}
\newcommand{\Z}{\mathbb{Z}}
\newcommand{\Q}{\mathbb{Q}}
\newcommand{\R}{\mathbb{R}}
\newcommand{\C}{\mathbb{C}}
\newcommand{\F}{\mathbb{F}}
\newcommand{\A}{\mathbb{A}}
\newcommand{\PP}{\mathbb{P}}
\newcommand{\E}{\mathbb{E}}
\newcommand{\eps}{\varepsilon}
\newcommand{\dd}{\,\mathrm{d}}
\begin{document}
\section*{UnsolvedMath ID: 800231}
\textbf{Source problem code:} catalogue entry
\subsection{Definitions and mathematical statement}
Let \(n\geq 1\), and let \(D\subset\C^n\) be a pseudoconvex domain with real-analytic boundary. Here real-analytic boundary means that each boundary point has a neighborhood \(U\) with a real-valued real-analytic defining function \(\rho\) such that \(D\cap U=\{\rho<0\}\) and \(\dd\rho\neq 0\) on \(\partial D\cap U\). The source describes \(D\) as ``real analytically bounded''; whether this also requires \(D\) to be bounded is not explicit in the supplied record.

For a fixed degree \(1\leq q\leq n\), let \(L^2_{(0,q)}(D)\) denote the space of square-integrable \((0,q)\)-forms with respect to Euclidean volume. Let \(\bar\partial\) be the maximal distributional Cauchy--Riemann operator, and let \(\bar\partial^{\ast}\) be its Hilbert-space adjoint. The \(\bar\partial\)-Neumann Laplacian is
\[
\square_q=\bar\partial\bar\partial^{\ast}+\bar\partial^{\ast}\bar\partial,
\]
with domain
\[
\operatorname{Dom}(\square_q)=\left\{u\in\operatorname{Dom}(\bar\partial)\cap\operatorname{Dom}(\bar\partial^{\ast}):\ \bar\partial u\in\operatorname{Dom}(\bar\partial^{\ast}),\ \bar\partial^{\ast}u\in\operatorname{Dom}(\bar\partial)\right\}.
\]
These adjoint-domain requirements encode the \(\bar\partial\)-Neumann boundary conditions.

Write \(\mathcal A^\omega_{(0,q)}(\overline D)\) for forms whose coefficients, near every point of \(\overline D\), are restrictions of functions real analytic in the \(2n\) real coordinates on an open neighborhood of that point.

The source asks whether the \(\bar\partial\)-Neumann problem on every such \(D\) is globally real-analytic hypoelliptic up to the boundary. A standard precise operator interpretation is whether
\[
\forall u\in\operatorname{Dom}(\square_q),\qquad
\square_q u\in\mathcal A^\omega_{(0,q)}(\overline D)
\ \Longrightarrow\ 
u\in\mathcal A^\omega_{(0,q)}(\overline D).
\]
The supplied statement does not specify the form degree, its quantifier, or the intended solution normalization. Consequently, the displayed implication is an explicit candidate interpretation, not a resolution of those omissions. In particular, the record does not determine whether the intended assertion concerns this implication for all solutions or analytic regularity of a distinguished \(\bar\partial\)-Neumann solution operator.
\subsection{Short English statement}
Does a pseudoconvex domain with real-analytic boundary always have global real-analytic regularity up to the boundary for the \(\bar\partial\)-Neumann problem? The source leaves the form degree, the precise solution formulation, and any additional boundedness requirement unspecified. It attributes the question to D. S. Tartakoff.
\subsection{Sources}
\begin{itemize}
\item Local UnsolvedMath registry, permanent ID \(800231\), exact source problem code \texttt{catalogue entry}, title ``Several Complex Variables Problem List --- Problem 28.'' The registry labels this record open; that label is supplied source metadata, not an independently verified mathematical assessment.
\item Buzzard, Fornaess, and Krantz, \emph{The Several Complex Variables Problem List} (1995), Problem 28. Supplied source URL: \url{https://arxiv.org/abs/math/9509202}. The record reports access on 2026-10-16 and extraction by \texttt{arxiv\_latex}; these are supplied provenance statements and have not been independently verified here.
\item The problem text attributes the question to D. S. Tartakoff and supplies the historical contact \texttt{U22393@UICVM.bitnet}. The registry lists Buzzard, Fornaess, and Krantz in its proposer field and gives 1995 as the proposed year; the specific attribution within the problem text is retained separately.
\end{itemize}
\end{document}
